\paragraph*{Objectif du chapitre :}
\begin{itemize}
\item  Connaitre ce que sont $\mathbb{N}$ et $\mathbb{Z}$ ainsi que leurs noms et notations.
\item  Appliquer les règles de priorités opératoires.
\item  Faire les additions et multiplications de nombres négatifs ou positifs.
\item  Connaitre le nom du résultat d'un calcul.
\end{itemize}

\section{Entiers naturels : $\mathbb{N}$}

\begin{defi}{}{}
\noindent Les entiers naturels sont les nombres qui servent à dénombrer, à compter.

\noindent L'ensemble de tous les entiers naturels est: $$\mathbb{N}=\{ 0,\ 1,\ 2,\ 3,\ \dots \}$$

\noindent Ce sont des nombres positifs ou nuls.
\end{defi}

\begin{rem}{}{}
\begin{itemize}
\item La notation $\mathbb{N}$ représente la première lettre de \og naturel\fg{}.
\item En fait l'ensemble des nombres entiers naturels se construit logiquement.
\item Parmi les mathématiciens qui ont marqué l'étude de l'ensemble des entiers naturels : 

Hermann Günther Grassmann (XIX$^e$ siècle), Giuseppe Peano (fin XIX$^e$ siècle).

\medskip
\href{https://www.axelnax.fr/index.php?p=frise}{Voir la frise chronologique}
\end{itemize}
\end{rem}

\subsection{Divisibilité}

\begin{defi}{}{}
Un nombre entier $a$ est divisible par un nombre entier $b$ si le reste de la division euclidienne de $a$ par $b$ est nul. 

\noindent Dans ce cas, $a$ est un multiple de $b$ et $b$ est un diviseur de $a$.
\end{defi}

\smallskip
\begin{ex}{}{}
$72$ est divisible par $24$. En effet, $72 = 3\times 24 + 0$ donc le reste de la division euclidienne de $72$ par $24$ est nul.

\noindent On peut donc dire que $72$ est un multiple de $24$ et de $3$ ou que $24$ est un diviseur de $72$.
\end{ex}

\begin{defi}{}{}
Un entier multiple de deux est un entier pair, les autres sont les entiers impairs.
\end{defi}

\smallskip
\begin{ex}{}{}
$72$ est un multiple de $2$. En effet, $2\times 36 = 72$. Donc 72 est un entier pair.
\end{ex}

\begin{prop}{}{}
Soit $a$ un entier naturel, la somme de deux multiples de $a$ est multiple de $a$.
\end{prop}

\begin{deme}{}{}
Soient $b$ et $c$ deux nombres entiers.

\noindent On cherche à démonter l'implication suivante : 
\begin{center}
Si $b$ est un multiple de $a$ et $c$ un multiple de $a$, alors $b+c$ est un multiple de $a$.
\end{center}

\medskip
Si $b$ est un multiple de $a$ et $c$ un multiple de $a$ alors il existe $m$ et $m'$ deux entiers tel que :
$$ b =m\times a  \quad \text{ et } \quad c = m' \times a$$

Par addition, on a :
$$ b + c = m\times a  + m' \times a $$

En factorisant par $a$, il vient : 
$$ b + c = (m  + m') \times a $$

Donc $b+c$ est un multiple de $a$. \hfill $\blacksquare$
\end{deme}

\begin{prop}{}{}
Le carré d'un nombre impair est impair.
\end{prop}

\begin{deme}{}{}
Soit $a$ un nombre entier.

\noindent On cherche à démonter l'implication suivante : 
\begin{center}
Si $a$ est un nombre impair, alors $a^2$ est un nombre impair.
\end{center}

\medskip
Si $a$ est un nombre impair alors il existe $k$ entiers tel que :
$$ a =2\times k+1 $$

Par multiplication, on a :
$$ a^2 = \left(2\times k+1 \right)^2 = (2\times k)^2 + 2 \times 2k \times 1 + 1^2 = 4k^2 + 4k +1$$

$4k^2$ est un multiple de $2$ ($2 \times 2k^2$)

$4k$ est multiple de $2$ ($2 \times 2k$)

D'après la propriété précédente, $4k^2+4k$ est donc un multiple de deux, donc un nombre pair.

Donc $a^2 = (4k^2 + 4k) +1$ est un nombre qui est après un nombre pair donc un nombre impair. \hfill $\blacksquare$
\end{deme}

\subsection{Critère de divisibilité}

\begin{methode}{}{}
\begin{itemize}
\item  Un nombre est divisible par $2$ si il se termine par $0$, $2$, $4$, $6$ ou $8$
\item  Un nombre est divisible par $5$ s'il se termine par $0$ ou $5$.

\item  Un nombre est divisible par $3$ si la somme de ses chiffres est divisible par $3$.

Exemple : $249$ est divisible par $3$ car $2 + 4 + 9 = 15$ qui est divisible par $3$ ($3 \times 5 = 15$).

\item  Un nombre est divisible par $9$ si la somme de ses chiffres est divisible par $9$.

Exemple : $2 538$ est divisible par $9$ car $2 + 5 + 3 + 8 = 18$ qui est divisible par $9$ ($3 \times 9 = 18$).
\end{itemize}
\end{methode}

\begin{defi}{}{}
Un nombre qui n'est divisible que par $1$ est lui-même est appelé un nombre premier.
\end{defi}

\begin{rem}{}{}
Il n'existe pas à ce jour de formule pour trouver des nombres premiers. Afin d'être certains qu'un nombre soit premier, il faut faire de nombreux calculs, souvent a l'aide d'ordinateur.
\end{rem}

\newpage
\section{Entiers relatifs : $\mathbb{Z}$}

\begin{defi}{}{}
\noindent Les entiers relatifs sont les entiers naturels et leurs opposés.

\noindent L'ensemble de tous les entiers relatifs est:
$$\mathbb{Z}=\{ \dots ,\ -3,\ -2,\ -1,\ 0,\ 1,\ 2,\ 3,\ \dots \}$$
\end{defi}

\begin{prop}{}{}
Les nombres entiers relatifs ont une partie décimale nulle.
\end{prop}

\begin{rem}{}{}
\begin{itemize}
\item Autrement dit l'écriture décimale infinie d'un entier ne comporte que des zéros dans la partie décimale. 

Ainsi: $3=3,000\dots$
\item Il existe des nombres qui ne sont pas des entiers. Par exemple $\dfrac{1}{2}$ n'est pas un entier.
\item La notation $\mathbb{Z}$ (de l'allemand \og zahlen\fg{} signifiant nombres) fut popularisée par Nicolas Bourbaki.
\end{itemize}
\end{rem}

%\subsection{Additionner des entiers relatifs.}

%\begin{ex}{}{}
%Évaluons les quantités suivantes:

%\begin{tabularx}{\linewidth}{*{2}{X}}
%\begin{minipage}{\linewidth}
%\begin{itemize}
%\item $A=-3+2$
%\item $B=-8-7$
%\end{itemize}	
%\end{minipage}	
%&
%\begin{minipage}{\linewidth}			
%\begin{itemize}
%\setcounter{enumi}{2}
%\item $C=-(-3)+4$
%\item $D=-13+(-11)$
%\end{itemize}
%\end{minipage}
%\end{tabularx}
%\end{ex}

%\subsection{Multiplier des entiers relatifs.}

\begin{prop}{}{}
\noindent Le produit de deux nombres de même signe est positif.

\noindent Le produit de deux nombres de signes contraires est négatif.
\end{prop}

%\medskip
%\noindent Il est également possible de se représenter les choses comme suit.

%\bigskip
%\begin{minipage}{7cm}
%\begin{center}
%\begin{tikzpicture}
%\coordinate (A) at (0,0);
%\coordinate (B) at (4,0);
%\coordinate (C) at (2,1);	
%\node[draw,rectangle] (A) at (A) {$x$};
%\node[draw,rectangle] (B) at (B) {$ {\color{G1}-2 \times} x$};
%\draw[>=latex,->] (A)--(B)node[midway, above]{$\color{G1}-2 \times$};
%\draw node (C) at (C) {\footnotesize Multiplier par un nombre négatif change le signe};	
%\end{tikzpicture}
%\end{center}
%\end{minipage}
%\hspace{5mm}
%\begin{minipage}{7cm}
%\begin{center}
%\begin{tikzpicture}	
%\coordinate (A1) at (0,-2);
%\coordinate (B1) at (4,-2);
%\coordinate (C1) at (2,-1);	
%\node[draw,rectangle] (A1) at (A1) {$x$};
%\node[draw,rectangle] (B1) at (B1) {$ {\color{G1} 3 \times }x$};
%\draw[>=latex,->] (A1)--(B1)node[midway, above]{${\color{G1} 3 \times} $};
%\draw node (C1) at (C1) {\footnotesize Multiplier par un nombre positif ne change pas le signe};	
%\end{tikzpicture}
%\end{center}
%\end{minipage}

\begin{rem}{}{}
\noindent $0$ est dit \textbf{absorbant} pour la multiplication : \quad quelque soit le nombre $a$, $0\times a=0$.

\noindent $1$ est dit \textbf{neutre} pour la multiplication : \quad quelque soit le nombre $a$, $1\times a= a$.
\end{rem}

%\begin{ex}{}{}
%Évaluons les quantités suivantes:

%\begin{tabularx}{\linewidth}{*{2}{X}}
%\begin{minipage}{\linewidth}
%\begin{itemize}
%\item $A=-3\times 2$
%\item $B=-8\times (-7)$
%\end{itemize}	
%\end{minipage}	
%&
%\begin{minipage}{\linewidth}			
%\begin{itemize}
%\setcounter{enumi}{2}
%\item $C=-(-3)\times 4$
%\item $D=-13\times (-11)$
%\end{itemize}
%\end{minipage}
%\end{tabularx}
%\end{ex}

%\subsection{Priorités opératoires.}

\begin{prop}{}{}
Les règles de priorités pour l'instant se limitent à : (Dans l'ordre des priorités)
\begin{itemize}
	\item Les calculs entre parenthèses ou crochets. S'il y a des parenthèses emboîtées les plus emboîtées sont prioritaires.
	\item Les multiplications en allant de gauche à droite.
	\item Les additions et soustractions en allant de gauche à droite.
\end{itemize}
\end{prop}

%\begin{ex}{}{}

%Évaluons les quantités suivantes:
	
%\noindent 
%\begin{tabularx}{\linewidth}{*{2}{X}}
%\begin{minipage}{\linewidth}
%\begin{itemize}
%\item $A=4\times (-1)$
%\item $B=-1\times (-1)$
%\end{itemize}
%\end{minipage}
%&
%\begin{minipage}{\linewidth}
%\begin{itemize}
%\setcounter{enumi}{2}
%\item $C=3\times [-2(3-5)]$
%\item $D=(3-5\times (-2))(-1+(-4)\times 2)$
%\end{itemize}
%\end{minipage}
%\end{tabularx}
%\end{ex}

\begin{exec}
\hspace{0.5cm} \hypertarget{Ex_1_1_1}{} Exercice \ref{1_1_1} \hspace{0.5cm}
\end{exec}

\section{Nommer le résultat d'un calcul.}

\begin{defi}{}{}
Nous dirons qu'un calcul est une \textbf{somme} (respectivement une \textbf{différence}, respectivement un \textbf{produit}) si la dernière opération effectuée en respectant les priorités opératoires est une addition (respectivement une soustraction, respectivement une multiplication).
\end{defi}
\medskip

À moins de devoir distinguer somme et différence, les résultats d'une addition et d'une soustraction seront très souvent indifféremment appelés des sommes.

\begin{ex}{}{}
Ainsi \quad $2\times(-2) +4\times (3-1)$ \quad et \quad $3-(4-2)\times (5-2)$ \quad sont des sommes tandis que \quad $(3-1)\times (-6-1)$ \quad est un produit.
\end{ex}

\begin{exec}
\hspace{0.5cm} \hypertarget{Ex_1_1_3}{} Exercice \ref{1_1_3} \hspace{0.5cm};
\hspace{0.5cm} \hypertarget{Ex_1_1_4}{} Exercice \ref{1_1_4} \hspace{0.5cm}
\end{exec}